\documentclass[11pt,reqno]{amsart}
\usepackage{amsmath,amssymb,amsthm}
\usepackage[numbers]{natbib}
% \usepackage[dvips]{graphicx}
% \usepackage{epsfig,epstopdf}
% \usepackage[dvipdfm]{hyperref}

%\oddsidemargin +0.1in
% \evensidemargin +0.1in
%\topmargin 0pt \textheight 8.7in \textwidth 6.35in
%\linespread{1.3}\parskip .05in


%% ---------------------------------------------------------------
\vfuzz2pt % Don't report over-full v-boxes if over-edge is small
\hfuzz2pt % Don't report over-full h-boxes if over-edge is small
% THEOREMS -------------------------------------------------------
\newtheorem{theorem}{Theorem}
% \newtheorem{algorithm}{Algorithm}[section]
% \newtheorem{conjecture}{Conjecture}
\newtheorem{corollary}{Corollary}
\newtheorem{lemma}{Lemma}
\newtheorem{proposition}{Proposition}
% \newtheorem{assumption}{C\!\!}
\theoremstyle{definition}
\newtheorem{definition}{Definition}
% \newtheorem{step}{Step}
\newtheorem{remark}{Comment}
% \numberwithin{remark}{section}
\newtheorem{example}{Example}
%\numberwithin{equation}{section}


%\newcommand{\cite}{\cite}
% \newcommand{\cite}{\cite}

\newcommand{\eps}{\varepsilon}
% \newcommand{\eg}{e.g., \xspace}
% \newcommand{\ie}{i.e.,\xspace}
% \newcommand{\etc}{etc.\@\xspace}
% \newcommand{\iid}{\emph{i.i.d.}\ }
% \newcommand{\etal}{et.\ al.\ }
% \newcommand{\PP}{P}
% \newcommand{\be}{\begin{eqnarray}}
% \newcommand{\ee}{\end{eqnarray}}
% \newcommand{\ind}{ \overset{d}{\longrightarrow}}
% \newcommand{\inp}{ \overset{p}{\longrightarrow}}
% \newcommand{\ba}{\begin{array}}
% \newcommand{\ea}{\end{array}}
% \newcommand{\bs}{\begin{align}\begin{split}\nonumber}
% \newcommand{\bsnumber}{\begin{align}\begin{split}}
% \newcommand{\es}{\end{split}\end{align}}
% \newcommand{\doublespace}{\linespread{1.1}}
% \newcommand{\singlespace}{\linespread{0.85}}
% \newcommand{\fns}{\singlespace\footnotesize}
% \newcommand{\cooldate}{\begin{flushright}\footnotesize\today\normalsize\end{flushright}}

% \newcommand{\sss}{\scriptscriptstyle}
% \newcommand{\n}{n}
% \newcommand{\inpr}{ \overset{p^*_{\scriptscriptstyle n}}{\longrightarrow}}
% \newcommand{\T}{ \scriptscriptstyle T}
\newcommand{\LL}{ \mathcal{L}}
% \newcommand{\M}{ \scriptscriptstyle M}
% \newcommand{\al}{ \scriptscriptstyle A}
% \renewcommand{\(}{\left(}
% \renewcommand{\)}{\right)}
\renewcommand{\hat}{\widehat}

\newcommand{\Pn}{\Pr_n}
\newcommand{\Pp}{\mathbf{P}}
\newcommand{\Ep}{{\mathrm{E}}}
\newcommand{\barEp}{\bar \Ep}
\newcommand{\En}{{\mathbb{E}_n}}
% \newcommand{\Ena}{{\mathbb{E}_{n_a}}}
% \newcommand{\Enb}{{\mathbb{E}_{n_b}}}
% \newcommand{\Enk}{{\mathbb{E}_{n_k}}}
% \newcommand{\LASSO}{\sqrt{\text{Lasso}}}


\renewcommand{\Pr}{{\mathrm{P}}}
\newcommand{\barf}{\overline{f}}
% \newcommand{\barfp}{\overline{f'}}
% \newcommand{\Uniform}{{\text{Uniform}}}
% \newcommand{\diam}{{\text{diam}}}
% \newcommand{\trace}{{\text{trace}}}
\newcommand{\RR}{\mathbb{R}}
% \def\mK{{\mathcal K}}
% \newcommand{\expect}{\mathcal{E}}
% \def\z{{d}}
% \def\x{{x}}
% \def\supp{{\rm support}}
% \newcommand{\floor}[1]{\left\lfloor #1 \right\rfloor}
% \newcommand{\ceil}[1]{\left\lceil #1 \right\rceil}
% \newcommand{\semin}[1]{\phi_{{\rm min}}(#1)}
% \newcommand{\semax}[1]{\phi_{{\rm max}}(#1)}
% \renewcommand{\hat}{\widehat}
% \renewcommand{\leq}{\leqslant}
% \renewcommand{\geq}{\geqslant}
% \renewcommand{\qed}{\hfill{\tiny \ensuremath{\blacksquare} }}%
% \newcommand{\sign}{ {\rm sign}}

% \DeclareMathOperator{\pr}{pr}


\DeclareMathOperator{\Corr}{Corr}
\DeclareMathOperator{\Var}{Var}

\begin{document}\section{Proofs for Section \ref{sec: Gaus and multiplier}}
\label{sec: proof for section 2}

\subsection{Proof of Theorem  \ref{theorem:comparison}}

Here for a smooth function $f: \RR^{p} \to \RR$, we write $\partial_{j} f (z) = \partial f(z)/\partial z_{j}$ for $z = (z_{1},\dots,z_{p})^{T}$.
We shall use the following version of Stein's identity.

\begin{lemma}[Stein's identity]
\label{lemma: talagrand}
Let $W =(W_{1},\dots,W_{p})^{T}$ be a centered Gaussian random vector in $\RR^{p}$.
Let $f: \RR^{p} \to \RR$ be a $C^1$-function such that $\Ep [ |\partial_{j} f(W) | ] < \infty$ for all $1 \leq j \leq p$. Then for every $1 \leq j \leq p$,
\begin{equation*}
\Ep[W_{j}f(W)]=\sum_{k=1}^p\Ep[W_{j} W_{k}] \Ep [\partial_k f (W)].
\end{equation*}
\end{lemma}
\begin{proof}[Proof of Lemma \ref{lemma: talagrand}]
See Section A.6 of \cite{Talagrand2003}; also \cite{ChenGoldsteinShao2011} and \cite{Stein1981}.
\end{proof}

We will use the following properties of the smooth max function.

\begin{lemma}
\label{lemma: smooth max property}
For every $1 \leq j,k \leq p$,
\begin{equation*}
 \partial_{j} F_{\beta}(z) = \pi_{j}(z), \quad \partial_{j} \partial_k F_{\beta}(z) = \beta w_{jk} (z),
\end{equation*}
where
\begin{equation*}
\pi_{j}(z)  :=    e^{\beta z_{j}}/{\textstyle \sum}_{m=1}^pe^{\beta z_m}, \ w_{jk}(z) := 1 (j = k) \pi_{j} (z)- \pi_{j}(z) \pi_{k} (z).
\end{equation*}
% where $(z)$ denotes evaluation at $z$.
Moreover,
\begin{equation*}
\pi_{j}(z) \geq 0,  \  {\textstyle \sum}_{j=1}^p \pi_{j} (z) = 1, \ {\textstyle \sum}_{j,k=1}^p |w_{jk}(z)|  \leq 2.
\end{equation*}
\end{lemma}

\begin{proof}[Proof of Lemma \ref{lemma: smooth max property}]
The first property was noted in \cite{Chatterjee2005b}.  The other properties follow from a direct calculation.
\end{proof}

\begin{lemma}\label{lemma: second deriv m}
Let $m := g \circ F_{\beta}$ with $g \in C^{2}(\RR)$. Then for every $1 \leq j,k \leq p$,
\begin{align*}
\partial_{j} \partial_k m(z) =   (g'' \circ F_\beta) (z) \pi_{j} (z) \pi_{k} (z)+ \beta (g' \circ F_\beta) (z) w_{jk}(z),
\end{align*}
where $\pi_{j}$ and $w_{jk}$ are defined in Lemma \ref{lemma: smooth max property}.

\end{lemma}
\begin{proof}[Proof of lemma \ref{lemma: second deriv m}]
The proof follows from a direct calculation.
\end{proof}



\begin{proof}[Proof of Theorem \ref{theorem:comparison}]
Without loss of generality, we may assume that $X$ and $Y$ are independent, so that $\Ep[X_{j} Y_k]=0$ for all $1 \leq j,k \leq p$.
Consider the following Slepian interpolation between $X$ and $Y$:
\begin{equation*}
Z(t):=\sqrt{t} X+\sqrt{1-t} Y, \ t \in [0,1].
\end{equation*}
Let $m := g \circ F_\beta$ and $\Psi(t):=\Ep[m (Z(t))]$.
Then
\begin{equation*}
|\Ep[m(X)]-\Ep[ m(Y) ]|= | \Psi (1) - \Psi (0)| =\left |\int_{0}^{1} \Psi'(t) dt\right|.
\end{equation*}
Here we have
\begin{align*}
\Psi'(t) & = \frac{1}{2} \sum_{j=1}^p \Ep [ \partial_{j} m(Z(t)) (t^{-1/2}X_{j}-(1-t)^{-1/2}Y_{j}) ]  \\
& = \frac{1}{2} \sum_{j=1}^p \sum_{k=1}^p (\sigma_{jk}^{X}-\sigma_{jk}^{Y}) \Ep[\partial_{j} \partial_k m (Z(t)) ],
\end{align*}
where the second equality follows from applying Lemma \ref{lemma: talagrand} to $W = (t^{-1/2}X_{j}-(1-t)^{-1/2}Y_{j}, Z(t)^{T})^{T}$
and $f(W)=\partial_{j} m(Z(t))$. Hence
\begin{align*}
\left |\int_{0}^{1} \Psi'(t) dt\right|
&\leq \frac{1}{2} \sum_{j,k=1}^{p} | \sigma_{jk}^{X}-\sigma_{jk}^{Y} | \left | \int_0^1 \Ep [\partial_{j} \partial_k m(Z(t))] dt \right| \\
&\leq \frac{1}{2} \max_{1 \leq j,k \leq p} | \sigma_{jk}^{X}-\sigma_{jk}^{Y} | \int_0^1 \sum_{j,k=1}^{p} \left | \Ep[\partial_{j} \partial_k m(Z(t))]\right | dt   \\
&= \frac{\Delta}{2} \int_0^1 \sum_{j,k=1}^{p} \left | \Ep[\partial_{j} \partial_k m(Z(t))]\right | dt.
\end{align*}
By Lemmas \ref{lemma: smooth max property} and \ref{lemma: second deriv m},
\begin{equation*}
\sum_{j,k=1}^{p} |\partial_{j} \partial_k m(Z(t))|   \leq   | (g'' \circ F_{\beta})(Z(t))| + 2 \beta | (g' \circ F_{\beta})(Z(t))|.
\end{equation*}
Therefore, we have
\begin{align*}
&\left|\Ep[g(F_{\beta}(X))-g(F_{\beta}(Y))]\right| \\
&\leq  \Delta \times  \left \{  \frac{1}{2} \int_0^1 \Ep[ | (g'' \circ F_{\beta})(Z(t))| ] dt +  \beta \int_0^1 \Ep [| (g' \circ F_{\beta})(Z(t))| ]dt\right \}  \\
&\leq  \Delta ( \| g'' \|_{\infty}/2 + \beta \| g' \|_{\infty}),
\end{align*}
which leads to the first assertion. The second assertion follows from the inequality (\ref{eq: smooth max property}). This completes the proof.
\end{proof}

\subsection{Proof of Theorem \ref{cor: distances Gaussian to Gaussian}}
The final assertion follows from the inequality $a_{p} \leq \sqrt{2 \log p}$ (see, for example, \cite{Talagrand2003}, Proposition 1.1.3).
Hence we prove (\ref{eq: K-distance}).
We first note that we may assume that $0 < \Delta < 1$ since otherwise the proof is trivial (take $C \geq 2$ in (\ref{eq: K-distance})).
In what follows, let  $C>0$ be a generic constant that depends only on $\min_{1 \leq j \leq p} \sigma_{jj}^{Y}$ and $\max_{1 \leq j \leq p} \sigma_{jj}^{Y}$, and its value may change from place to place.
For $\beta > 0$, define $e_{p,\beta} := \beta^{-1} \log p$.
Consider and fix a $C^{2}$-function $g_0:\RR \to [0,1]$ such that $g_0(t)=1$ for $t \leq 0$ and $g_0(t)=0$ for $t \geq 1$. For example, we may take
\begin{equation*}
g_{0}(t) =
\begin{cases}
0, & t \geq 1, \\
30\int_{t}^{1} s^{2}(1-s)^{2} ds, & 0 < t < 1, \\
1, & t \leq 0.
\end{cases}
\end{equation*}
For given $x \in \RR, \beta > 0$ and $\delta > 0$,  define $g_{x,\beta,\delta} (t)=g_0(\delta^{-1}(t-x-e_{p,\beta}))$.  For this function $g_{x,\beta,\delta}$, $\| g_{x,\beta,\delta}' \|_{\infty} = \delta^{-1} \| g_{0}' \|_{\infty}$ and $\| g_{x,\beta,\delta}'' \|_{\infty} = \delta^{-2} \| g_{0}'' \|_{\infty}$. Moreover,
\begin{equation}
1(t \leq x+e_{p,\beta}) \leq g_{x,\beta,\delta}(t) \leq 1(t \leq x + e_{p,\beta} + \delta), \ \forall t \in \RR. \label{eq: property of g}
\end{equation}

For arbitrary $x \in \RR, \beta > 0$ and $\delta > 0$, observe that
\begin{align}
\Pr ( \max_{1 \leq j \leq p} X_{j}   \leq x )
&\leq \Pr ( F_{\beta} (X) \leq x + e_{p,\beta}) \leq \Ep[g_{x,\beta,\delta}(F_\beta(X))] \notag \\
&\leq \Ep[g_{x,\beta,\delta}(F_\beta(Y))] + C(\delta^{-2}+\beta \delta^{-1})\Delta \notag  \\
&\leq  \Pr ( F_\beta (Y) \leq x+e_{p,\beta}+\delta ) + C(\delta^{-2}+\beta \delta^{-1})\Delta \notag \\
&\leq \Pr (\max_{1 \leq j \leq p} Y_{j} \leq x+e_{p,\beta}+\delta ) + C(\delta^{-2}+\beta \delta^{-1})\Delta, \label{eq: Kdistance-intermediate}
\end{align}
where the first inequality follows from the inequality (\ref{eq: smooth max property}), the second from the inequality (\ref{eq: property of g}),  the third from Theorem \ref{theorem:comparison}, the fourth from the inequality (\ref{eq: property of g}), and the last from the inequality (\ref{eq: smooth max property}).
We wish to compare $\Pr (\max_{1 \leq j \leq p} Y_{j} \leq x+e_{p,\beta}+\delta) $ with $\Pr (\max_{1 \leq j \leq p} Y_{j} \leq x ) $, and this is where the anti-concentration inequality plays its role.
By Theorem \ref{thm: anticoncentration}, we have
\begin{align}
&\Pr (\max_{1 \leq j \leq p} Y_{j} \leq x+e_{p,\beta}+\delta ) -  \Pr (\max_{1 \leq j \leq p} Y_{j} \leq x ) \notag \\
&= \Pr (x < \max_{1 \leq j \leq p} Y_{j} \leq x+e_{p,\beta}+\delta ) \leq \LL ( \max_{1 \leq j \leq p} Y_{j}, e_{p,\beta}+\delta) \label{eq: Kdistance-intermediate2} \\
&\leq  C (e_{p,\beta} + \delta) \sqrt{1 \vee a_{p}^{2} \vee \log \{ 1 / (e_{p,\beta} + \delta) \}} \leq  C (e_{p,\beta} + \delta) \sqrt{1 \vee a_{p}^{2} \vee  \log (1 / \delta)}. \notag
\end{align}
Therefore,
\begin{multline}
\Pr ( \max_{1\leq j\leq p}X_{j}\leq x ) -  \Pr (\max_{1\leq j\leq p}Y_{j}\leq x)  \\
\leq C  \left \{ (\delta^{-2}+\beta \delta^{-1})\Delta +  (e_{p,\beta}+\delta)\sqrt{1 \vee a_{p}^{2} \vee \log(1/\delta)}  \right \}. \label{normalcomp}
\end{multline}
Take $a = a_{p} \vee \log^{1/2}(1/\Delta)$, and choose  $\beta$ and $\delta$ in such a way that
\begin{equation*}
\beta = \delta^{-1}\log p \ \text{and} \ \delta = \Delta^{1/3} (1 \vee a)^{-1/3} (2\log p)^{1/3}.
\end{equation*}
Recall that $p \geq 2$ and $0 < \Delta < 1$. Observe that $(\delta^{-2}+\beta \delta^{-1})\Delta \leq C \Delta^{1/3} (1 \vee a)^{2/3} \log^{1/3} p, (e_{p,\beta}+\delta) (1 \vee a_{p}) \leq C \Delta^{1/3} (1 \vee a)^{2/3} \log^{1/3} p$, and since $\delta \geq \Delta^{1/3} (1 \vee a)^{-1/3}$, we have  $\log (1/\delta) \leq (1/3) \log ((1 \vee a)/\Delta)$. Hence the right side on (\ref{normalcomp}) is bounded by $C \Delta^{1/3} \{ (1 \vee a)^{2/3} \log^{1/3} p + (1 \vee a)^{-1/3} ( \log^{1/3} p ) \log^{1/2}((1 \vee a)/\Delta) \}$. In addition, $(1 \vee a)^{-1} \log^{1/2} (1 \vee a)$ is bounded by a universal constant, so that
\begin{multline*}
\text{right side of (\ref{normalcomp})} \\
\leq C \Delta^{1/3} \{ (1 \vee a)^{2/3} \log^{1/3} p + (1 \vee a)^{-1/3} ( \log^{1/3} p ) \log^{1/2}(1/\Delta) \}.
\end{multline*}
The second term inside the bracket is bounded by $( \log^{1/3} p ) \log^{1/3}(1/\Delta)$ as $(1 \vee a)^{-1/3} \leq (\log (1/\Delta))^{-1/6}$,
and first term is bound by $(1 \vee a_{p})^{2/3} \log^{1/3} p + ( \log^{1/3} p ) \log^{1/3}(1/\Delta)$. Adjusting the constant $C$, the right side on the above displayed equation is bounded by $C \Delta^{1/3} \{ 1 \vee a_{p}^{2} \vee \log(1/\Delta) \}^{1/3} \log^{1/3}p$.

For the opposite direction, observe that
\begin{align*}
\Pr ( \max_{1 \leq j \leq p} X_{j}   \leq x )
&\geq \Pr ( F_{\beta} (X) \leq x) \geq \Ep[g_{x-e_{p,\beta}-\delta,\beta,\delta}(F_\beta(X))] \\
&\geq \Ep[g_{x-e_{p,\beta}-\delta,\beta,\delta}(F_\beta(Y))] - C(\delta^{-2}+\beta \delta^{-1})\Delta \\
&\geq \Pr (F_{\beta}(Y) \leq x - \delta ) - C(\delta^{-2}+\beta \delta^{-1})\Delta \\
&\geq \Pr (\max_{1 \leq j \leq p} Y_{j} \leq x - e_{p,\beta} -  \delta ) - C(\delta^{-2}+\beta \delta^{-1})\Delta.
\end{align*}
The rest of the proof is similar and hence omitted.
\qed


\section{Proof of Theorem \ref{thm: anticoncentration}}
\label{sec: proof for section 3}

The proof of Theorem  \ref{thm: anticoncentration} uses some properties of Gaussian measures.
We begin with preparing technical tools. The following two facts were essentially noted in \cite{Y65,Y68} (note: \cite{Y65} and \cite{Y68} did not contain a proof of Lemma \ref{lemma: Y65-1}, which we find non-trivial). For the sake of completeness, we give their proofs after the proof of Theorem \ref{thm: anticoncentration}.
\begin{lemma}
\label{lemma: Y65-1}
Let $(W_{1},\dots,W_{p})^{T}$ be a (not necessarily centered) Gaussian random vector in $\RR^{p}$ with $\Var(W_{j}) = 1$ for all $1 \leq j \leq p$.  Suppose that $\Corr(W_{j},W_{k})<1$ whenever $j \neq k$. Then the distribution of $\max_{1 \leq j \leq p} W_{j}$ is absolutely continuous with respect to the Lebesgue measure and a version of the density is given by
\begin{equation}
f(x) = \phi(x) \sum_{j=1}^{p} e^{ \Ep[W_{j}] x-(\Ep[W_{j}])^2/2 } \cdot \Pr\left (W_{k} \leq x, \forall k \neq j \mid W_{j} = x \right). \label{density}
\end{equation}
\end{lemma}

\begin{lemma}
\label{lemma: Y65-2}
Let $(W_{0},W_{1},\dots,W_{p})^{T}$ be  a (not necessarily centered) Gaussian random vector with $\Var(W_{j}) =1 $ for all $0 \leq j \leq p$. Suppose that $\Ep[W_{0} ]\geq 0$. Then the map
\begin{equation}\label{eq: prob}
x \mapsto  e^{ \Ep[W_{0} ] x-(\Ep[W_{0}])^2/2 } \cdot\Pr( W_{j} \leq x, 1 \leq \forall j \leq p  \mid W_{0} = x)
\end{equation}
is non-decreasing on $\RR$.
\end{lemma}

Let us also recall (a version of) the Gaussian concentration (more precisely, deviation) inequality. See,  for example,  \cite{L01}, Theorem 7.1, for its proof.
\begin{lemma}\label{thm3}
Let $(X_{1},\dots,X_{p})^{T}$ be a centered Gaussian random vector in $\RR^{p}$ with $\max_{1 \leq j \leq p} \Ep [X_{j}^{2}] \leq  \sigma^{2}$ for some $\sigma^{2}> 0$. Then for every $r > 0$,
\begin{equation*}
\Pr  ( \max_{1 \leq j \leq p} X_{j} \geq \Ep [ \max_{ 1 \leq j \leq p} X_{j}  ] + r ) \leq e^{-r^{2}/(2\sigma^{2})}.
\end{equation*}
\end{lemma}

We are now in position to prove Theorem \ref{thm: anticoncentration}.

\begin{proof}[Proof of Theorem \ref{thm: anticoncentration}]
The proof consists of three steps.

{\bf Step 1}. This step reduces the analysis to the unit variance case. Pick any $x\geq0$. Let $W_{j} :=(X_{j}-x)/\sigma_{j}+x/\underline{\sigma}$.
Then $\Ep[W_{j}]\geq 0$ and $\Var (W_{j})=1$. Define $Z:= \max_{1 \leq j \leq p} W_{j}$.
Then we have
\begin{align*}
&\Pr (  |\max_{1 \leq j \leq p}X_{j}-x |\leq \epsilon ) \leq \Pr\left(\left|\max_{1 \leq j \leq p}\frac{X_{j}-x}{\sigma_j}\right|\leq \frac{\epsilon}{\underline{\sigma}}\right) \\
&\quad \leq \sup_{y\in\RR}\Pr\left(\left|\max_{1 \leq j \leq p}\frac{X_{j}-x}{\sigma_j}+\frac{x}{\underline{\sigma}}-y\right|\leq \frac{\epsilon}{\underline{\sigma}}\right) =\sup_{y\in\RR}\Pr\left(\left|Z-y\right|\leq \frac{\epsilon}{\underline{\sigma}}\right).
\end{align*}


{\bf Step 2}. This step bounds the density of $Z$.
Without loss of generality, we may assume  that $\Corr (W_{j},W_{k})<1$ whenever $j \neq k$. Since the marginal distribution of $W_{j}$ is $N(\mu_{j},1)$ where $\mu_{j}:=\Ep[W_{j}]=(x/\underline{\sigma}-x/\sigma_j)\geq 0$, by Lemma \ref{lemma: Y65-1}, $Z$ has density of the form
\begin{equation}\label{eq: density form}
f_{p}(z) = \phi(z) G_{p} (z),
\end{equation}
where the map $z \mapsto G_{p}(z)$ is non-decreasing by Lemma \ref{lemma: Y65-2}. Define $\bar{z}:= (1/\underline{\sigma}-1/\overline{\sigma})x$, so that $\mu_{j}\leq \bar{z}$ for every $1 \leq j \leq p$. Moreover, define $\bar{Z}:=\max_{1 \leq j \leq p}(W_{j}-\mu_{j})$.  Then
\begin{align*}
&\int_{z}^{\infty} \phi (u) du G_{p}(z) \leq \int_{z}^{\infty} \phi(u) G_{p}(u) du = \Pr( Z> z) \\
&\qquad \leq \Pr( \bar{Z} > z - \bar{z} ) \leq  \exp \left \{ -\frac{(z-\bar{z}-\Ep[ \bar{Z}  ])_{+}^{2}}{2} \right \},
\end{align*}
where the last inequality is due to the Gaussian concentration inequality (Lemma \ref{thm3}).
Note that $W_{j}-\mu_{j}=X_{j}/\sigma_{j}$, so that
\begin{equation*}
\Ep[\bar{Z}]=\Ep [ \max_{1 \leq j \leq p}(X_{j}/\sigma_{j}) ] =: a_p.
\end{equation*}
Therefore, for every $z \in \RR$,
\begin{equation}\label{eq: bound Gn}
G_{p} (z) \leq \frac{1}{1-\Phi(z)} \exp \left \{ -\frac{(z-\bar{z}-a_p)_{+}^{2}}{2} \right \}.
\end{equation}
Mill's inequality states that for $z>0$,
\begin{equation*}
z \leq \frac{\phi(z)}{1-\Phi (z)}\leq z \frac{1+z^2}{z^2},
\end{equation*}
and in particular $(1+z^2)/z^2 \leq 2$ when $z>1$. Moreover, $\phi(z)/\{ 1-\Phi(z) \}\leq 1.53 \leq 2 $ for  $z \in (-\infty, 1)$.  Therefore,
\begin{equation*}
\phi(z)/\{ 1-\Phi(z) \} \leq 2 (z \vee 1),  \ \forall z \in \RR.
\end{equation*}
Hence we conclude from this, (\ref{eq: bound Gn}), and (\ref{eq: density form}) that
\begin{equation*}
f_{p}(z) \leq 2(z \vee 1) \exp \left \{  -\frac{(z-\bar{z}-a_p)_{+}^{2}}{2} \right \}, \ \forall z \in \RR.
\end{equation*}


{\bf Step 3}. By Step 2, for every $y \in \RR$ and $t > 0$,  we have
\begin{align*}
\Pr\left( | Z - y | \leq t \right) &= \int_{y -t}^{y+t} f_{p}(z) dz
\leq 2 t \max_{z \in [y -t, y+t]}f_{p}(z)  \leq 4 t (\bar{z}+ a_p + 1),
\end{align*}
where the last inequality follows from the fact that the map $z \mapsto z e^{-(z-a)^{2}/2}$ (with $a > 0$) is non-increasing on $[a+1,\infty)$.
Combining this inequality with  Step 1, for every $x \geq 0$ and $\epsilon > 0$, we have
\begin{equation}
\Pr ( | \max_{1 \leq j \leq p}X_{j} - x | \leq \epsilon ) \leq 4 \epsilon \{  (1/\underline{\sigma}-1/\overline{\sigma}) |x|+a_p + 1 \} /\underline{\sigma}.  \label{eq: anticoncentration main}
\end{equation}
This inequality also holds for $x<0$ by the similar argument, and hence it holds
for every $x \in \RR$.

If  $\underline{\sigma}=\overline{\sigma} =\sigma$, then we have
\begin{equation*}
\Pr ( | \max_{1 \leq j \leq p}X_{j} - x | \leq \epsilon ) \leq 4 \epsilon (a_p + 1)/\sigma, \ \forall  x \in \RR, \ \forall \epsilon > 0,
\end{equation*}
which leads to the first assertion of the theorem.

On the other hand, consider the case where $\underline{\sigma} < \overline{\sigma}$. Suppose first that $0 < \epsilon \leq \underline{\sigma}$.
By the Gaussian concentration inequality (Lemma \ref{thm3}), for $|x|\geq \epsilon+\overline \sigma(a_p+\sqrt{2\log(\underline{\sigma}/\epsilon)})$, we have
\begin{align}
&\Pr ( | \max_{1 \leq j \leq p} X_{j}-x|\leq\epsilon ) \leq \Pr (\max_{1 \leq j \leq p}X_{j} \geq |x|-\epsilon ) \notag \\
&\quad \leq \Pr \left (\max_{1 \leq j \leq p}X_{j}\geq \Ep[\max_{1 \leq j \leq p}X_{j}] +\overline{\sigma}\sqrt{2\log(\underline{\sigma}/\epsilon)} \right ) \leq \epsilon/\underline{\sigma}. \label{eq: anticoncentration tail}
\end{align}
For  $|x| \leq \epsilon+\overline \sigma(a_p+\sqrt{2\log(\underline{\sigma}/\epsilon)})$, by  (\ref{eq: anticoncentration main}) and using $\epsilon \leq \underline{\sigma}$, we have
\begin{align}
&\Pr ( | \max_{1 \leq j \leq p}X_{j} - x | \leq \epsilon ) \notag \\
&\quad \leq 4 \epsilon \{ (\overline{\sigma}/\underline{\sigma}) a_{p} +  (\overline{\sigma}/\underline{\sigma}-1)\sqrt{2\log(\underline{\sigma}/\epsilon)}   + 2 - \underline{\sigma}/\overline{\sigma}\}/\underline{\sigma}.
 \label{eq: anticoncentration non-tail}
\end{align}
Combining (\ref{eq: anticoncentration tail}) and (\ref{eq: anticoncentration non-tail}), we obtain the inequality in (ii) for $0 < \epsilon \leq \underline{\sigma}$ (with a suitable choice of $C$). If $\epsilon> \underline{\sigma}$, the inequality in (ii) trivially follows by taking $C \geq 1/\underline{\sigma}$.  This completes the proof.
\end{proof}


\begin{proof}[Proof of Lemma \ref{lemma: Y65-1}]
Let $M := \max_{1 \leq j \leq p} W_{j}$.
The absolute continuity of the distribution of $M$ is deduced from the fact that $\Pr(M \in A) \leq \sum_{j=1}^{p} \Pr(W_{j} \in A)$ for every Borel measurable subset $A$ of $\RR$.
Hence, to show that a version of the density of $M$ is given by (\ref{density}), it is enough to show that $\lim_{\epsilon \downarrow 0} \epsilon^{-1} \Pr(x < M \leq x+\epsilon)$ equals the right side on (\ref{density}) for a.e. $x \in \RR$.

For every $x \in \RR$ and $\epsilon > 0$, observe that
\begin{align*}
& \{ x < M \leq x + \epsilon \} \\
& = \{ \exists i_{0}, W_{i_{0}} > x \ \text{and} \ \forall i,  W_{i} \leq x +\epsilon  \} \\
&= \{ \exists i_{1}, x < W_{i_{1}} \leq x + \epsilon \ \text{and} \ \forall i \neq i_{1}, W_{i} \leq x \} \\
&\quad \cup \{ \exists i_{1}, \exists i_{2}, x < W_{i_{1}} \leq x + \epsilon, x < W_{i_{2}} \leq x + \epsilon \ \text{and} \ \forall i \notin \{ i_{1},i_{2} \},   W_{i} \leq x \} \\
&\qquad \qquad \vdots \\
&\quad \cup \{ \forall i, x < W_{i} \leq x + \epsilon \} \\
&=: A^{x,\epsilon}_{1} \cup A^{x,\epsilon}_{2} \cup \cdots \cup A^{x,\epsilon}_{p}.
\end{align*}
Note that the events $A^{x,\epsilon}_{1},A^{x,\epsilon}_{2},\dots,A^{x,\epsilon}_{p}$ are disjoint. For $A^{x,\epsilon}_{1}$, since
\begin{equation*}
A^{x,\epsilon}_{1} = \cup_{i=1}^{p} \{ x < W_{i} \leq x+\epsilon \ \text{and} \ W_{j} \leq x, \forall j \neq i \},
\end{equation*}
where the events on the right side are disjoint, we have
\begin{align*}
\Pr(A^{x,\epsilon}_{1}) &= \sum_{i=1}^{p} \Pr(x < W_{i} \leq x+\epsilon \ \text{and} \ W_{j} \leq x, \forall j \neq i) \\
&= \sum_{i=1}^{p} \int_{x}^{x+\epsilon} \Pr(W_{j} \leq x, \forall j \neq i \mid W_{i} = u) \phi(u-\mu_{i}) du,
\end{align*}
where $\mu_{i} := \Ep [ W_{i} ]$.
We show that for every $1 \leq i \leq p$ and a.e. $x \in \RR$, the map $u \mapsto \Pr(W_{j} \leq x, \forall j \neq i \mid W_{i} = u)$ is right continuous at $x$.
Let $X_{j}=W_{j}-\mu_{j}$ so that $X_{j}$ are standard Gaussian random variables. Then
\begin{equation*}
\Pr(W_{j} \leq x, \forall j \neq i \mid W_{i} = u)=\Pr(X_{j} \leq x-\mu_{j}, \forall j \neq i \mid X_{i} = u-\mu_{i}).
\end{equation*}
Pick $i=1$. Let $V_{j} = X_{j} - \Ep[ X_{j} X_{1}] X_{1}$ be the residual from the orthogonal projection of $X_{j}$ on $X_{1}$.
Note that the vector $(V_{j})_{2 \leq j \leq p}$ and $X_{1}$ are jointly Gaussian and uncorrelated, and hence independent, by which we have
\begin{align*}
&\Pr(X_{j} \leq x-\mu_{j}, 2 \leq \forall j \leq p \mid X_{1} = u-\mu_{1}) \\
&= \Pr(V_{j} \leq x-\mu_{j}-\Ep[ X_{j} X_{1}] (u-\mu_{1}), 2 \leq \forall j \leq p \mid X_{1} = u-\mu_{1}) \\
&=  \Pr(V_{j} \leq x-\mu_{j}-\Ep[ X_{j} X_{1}] (u-\mu_{1}), 2 \leq \forall j \leq p).
\end{align*}
Define $J := \{ j \in \{ 2,\dots,p \} : \Ep[ X_{j} X_{1}] \leq 0 \}$ and $J^{c} := \{ 2,\dots, p\} \backslash J$. Then
\begin{align*}
&\Pr(V_{j} \leq x-\mu_{j}-\Ep[ X_{j} X_{1}] (u-\mu_{1}), 2 \leq \forall j \leq p) \\
&\to \Pr( V_{j} \leq x_j, \forall j \in J, V_{j'} < x_{j'}, \forall j' \in J^{c}), \ \text{as} \ u \downarrow x,
\end{align*}
where $x_j=x-\mu_{j}-\Ep[X_{j}X_1](x-\mu_{1})$.
Here each $V_{j}$ either degenerates to $0$ (which occurs only when $X_{j}$ and $X_{1}$ are perfectly negatively correlated, that is, $\Ep [ X_{j} X_{1} ] = -1$) or has a non-degenerate Gaussian distribution, and hence for every $x \in \RR$ expect for at most $(p-1)$ points $(\mu_{1}+\mu_{j})/2, 2 \leq j \leq p$,
\begin{align*}
\Pr( V_{j} \leq x_j, \forall j \in J, V_{j'} < x_{j'}, \forall j' \in J^{c}) &=  \Pr( V_{j} \leq x_{j}, 2 \leq \forall j \leq p ) \\
&= \Pr(W_{j} \leq x, 2 \leq \forall j \leq p \mid W_{1} = x).
\end{align*}
Hence for $i=1$ and a.e. $x \in \RR$, the map $u \mapsto \Pr(W_{i} \leq x, \forall j \neq i \mid W_{i} = u)$ is right continuous at $x$. The same conclusion clearly holds for $2 \leq i \leq p$.
Therefore, we conclude that, for a.e. $x \in \RR$, as $\epsilon \downarrow 0$,
\begin{eqnarray*}
\frac{1}{\epsilon} \Pr(A_{1}^{x,\epsilon}) &\to& \sum_{i=1}^{p} \Pr(W_{j} \leq x, \forall j \neq i \mid W_{i} = x)\phi(x-\mu_{i})\\
&=&\phi(x)\sum_{i=1}^{p} e^{\mu_{i}x-\mu_{i}^2/2} \Pr(W_{j} \leq x, \forall j \neq i \mid W_{i} = x).
\end{eqnarray*}



In the rest of the proof, we  show that, for every $2 \leq i \leq p$ and $x \in \RR$, $\Pr(A_{i}^{x,\epsilon}) = o(\epsilon)$ as $\epsilon \downarrow 0$, which leads to the desired conclusion. Fix any $2 \leq i \leq p$. The probability $\Pr(A_{i}^{x,\epsilon})$ is bounded by a sum of terms of the form  $\Pr(x < W_{j} \leq x+\epsilon, x < W_{k} \leq x+\epsilon)$ with $j \neq k$. Recall that $\Corr(W_{j},W_{k})<1$. Assume that $\Corr(W_{j},W_{k})=-1$. Then for  every $x \in \RR$, $\Pr(x < W_{j} \leq x+\epsilon, x < W_{k} \leq x+\epsilon)$ is zero for sufficiently small $\epsilon$. Otherwise, $(W_{j},W_{k})^{T}$ obeys a two-dimensional, non-degenerate Gaussian distribution and hence $\Pr(x < W_{j} \leq x+\epsilon, x < W_{k} \leq x+\epsilon) = O(\epsilon^{2}) = o(\epsilon)$ as $\epsilon \downarrow 0$ for every $x \in \RR$. This completes the proof.
\end{proof}

\begin{proof}[Proof of Lemma \ref{lemma: Y65-2}]
Since $\Ep [ W_{0} ] \geq 0$, the map $x \mapsto \exp( \Ep[ W_{0} ] x- (\Ep[ W_{0} ])^{2} )$ is non-decreasing. Thus it suffices to show that the map
\begin{equation}\label{eq: reduced probability}
x \mapsto \Pr ( W_{1} \leq x, \dots,  W_{p} \leq x \mid W_{0} =x)
\end{equation}
is non-decreasing.
As in the proof of Lemma \ref{lemma: Y65-1}, let $X_{j}=W_{j}-\Ep [ W_{j} ]$ and let $V_{j} = X_{j} - \Ep[ X_{j} X_{0}] X_{0}$ be the residual from the orthogonal projection of $X_{j}$ on $X_{0}$.  Note that the vector $(V_{j})_{1 \leq j \leq p}$ and $X_{0}$ are independent.
Hence the  probability in (\ref{eq: reduced probability}) equals
\begin{align*}
&\Pr(V_{j} \leq x-\mu_{j} - \Ep[X_{j}X_{0}](x-\Ep[ W_{0} ] ),  1  \leq \forall j \leq p \mid X_{0}=x-\Ep[ W_{0} ])\\
&=\Pr(V_{j} \leq x-\mu_{j} - \Ep[X_{j}X_{0}](x-\Ep[ W_{0} ]),  1  \leq \forall j \leq p),
\end{align*}
where the latter is non-decreasing in $x$ on $\RR$ since $\Ep[X_{j}X_{0}] \leq 1$.
\end{proof}

\appendix

\section{Proof of Lemma \ref{lem: Delta estimate}}

Lemma \ref{lem: Delta estimate} follows from the following maximal inequality and H\"{o}lder's inequality. Here we write $a \lesssim b$  if $a$ is smaller than or equal to $b$ up to a universal positive constant.

\begin{lemma}
\label{lem: maximal ineq}
Let $Z_{1},\dots,Z_{n}$ be independent random vectors in $\RR^{p}$ with $p \geq 2$. Define $M := \max_{1 \leq i \leq n} \max_{1 \leq j \leq p} | Z_{ij} |$ and $\sigma^{2} := \max_{1 \leq j \leq p} \sum_{i=1}^{n} \Ep [ Z_{ij}^{2} ]$.
Then
\begin{equation*}
\Ep [\max_{1 \leq j \leq p} | {\textstyle \sum}_{i=1}^{n} (Z_{ij} - \Ep [Z_{ij}]) | ] \lesssim (\sigma \sqrt{\log p} + \sqrt{\Ep [ M^{2} ]} \log p).
\end{equation*}
\end{lemma}

We shall use the following lemma.

\begin{lemma}
\label{lem: maximal ineq nonnegative}
Let $V_{1},\dots,V_{n}$ be independent random vectors in $\RR^{p}$ with $p \geq 2$ such that $V_{ij} \geq 0$ for all $1 \leq i \leq n$ and $1 \leq j \leq p$. Then
\begin{equation*}
\Ep [ \max_{1 \leq j \leq p} {\textstyle \sum}_{i=1}^{n} V_{ij} ] \lesssim \max_{1 \leq j \leq p} \Ep [ {\textstyle \sum}_{i=1}^{n} V_{ij} ] + \Ep [ \max_{1 \leq i \leq n} \max_{1 \leq j \leq p} V_{ij} ] \log p.
\end{equation*}
\end{lemma}

\begin{proof}[Proof of Lemma \ref{lem: maximal ineq nonnegative}]
We make use of the symmetrization technique. Let $\eps_{1},\dots,\eps_{n}$ be independent Rademacher random variables (that is, $\Pr (\eps_{i} = 1) = \Pr ( \eps_{i} = -1) = 1/2$) independent of $V_{1}^{n} := \{ V_{1},\dots,V_{n} \}$.
Then by the triangle inequality and  Lemma 2.3.1 in \cite{VW96},
\begin{align*}
I := \Ep [ \max_{1 \leq j \leq p} {\textstyle \sum}_{i=1}^{n} V_{ij} ] &\leq \max_{1 \leq j \leq p} \Ep [ {\textstyle \sum}_{i=1}^{n} V_{ij} ] + \Ep [ \max_{1 \leq j \leq p} | {\textstyle \sum}_{i=1}^{n} (V_{ij} - \Ep [ V_{ij} ]) |] \\
&\leq \max_{1 \leq j \leq p} \Ep [ {\textstyle \sum}_{i=1}^{n} V_{ij} ] + 2\Ep [ \max_{1 \leq j \leq p} | {\textstyle \sum}_{i=1}^{n} \eps_{i} V_{ij}  |].
\end{align*}
By Lemmas 2.2.2 and 2.2.7 in \cite{VW96}, we have
\begin{align*}
\Ep [ \max_{1 \leq j \leq p} | {\textstyle \sum}_{i=1}^{n} \eps_{i} V_{ij}  | \mid V_{1}^{n} ] &\lesssim \max_{1 \leq j \leq p} ({\textstyle \sum}_{i=1}^{n} V_{ij}^{2})^{1/2}  \sqrt{\log p}\\
&\leq   \sqrt{B \log p} \max_{1 \leq j \leq p} ({\textstyle \sum}_{i=1}^{n} V_{ij})^{1/2},
\end{align*}
where $B := \max_{1 \leq i \leq n} \max_{1 \leq j \leq p} V_{ij}$. Hence by Fubini's theorem and the Cauchy-Schwarz inequality,
\begin{align*}
\Ep [ \max_{1 \leq j \leq p} | {\textstyle \sum}_{i=1}^{n} \eps_{i} V_{ij}  | ] &\lesssim \sqrt{\Ep [ B ] \log p} ( \Ep [ \max_{1 \leq j \leq p} {\textstyle \sum}_{i=1}^{n} V_{ij} ])^{1/2} \\
& = \sqrt{\Ep [ B ] \log p}  \sqrt{I}.
\end{align*}
Therefore, we have
\begin{equation*}
I \lesssim \max_{1 \leq j \leq p} \Ep [ {\textstyle \sum}_{i=1}^{n} V_{ij} ] + \sqrt{\Ep [ B ] \log p} \sqrt{I} =: a + b \sqrt{I}.
\end{equation*}
Solving this inequality, we conclude that $I \lesssim a + b^{2}$.
\end{proof}

\begin{proof}[Proof of Lemma \ref{lem: maximal ineq}]
Let $\eps_{1},\dots,\eps_{n}$ be independent Rademacher random variables independent of $Z_{1},\dots,Z_{n}$. Then arguing as in the previous proof, we have
\begin{align*}
\Ep [\max_{1 \leq j \leq p} | {\textstyle \sum}_{i=1}^{n} (Z_{ij} - \Ep [Z_{ij}]) | ]
&\leq 2 \Ep [\max_{1 \leq j \leq p} | {\textstyle \sum}_{i=1}^{n} \eps_{i}Z_{ij} | ] \\
&\lesssim \Ep [ \max_{1 \leq j \leq p} ( {\textstyle \sum}_{i=1}^{n} Z_{ij}^{2} )^{1/2} ]  \sqrt{\log p} \\
&\leq (\Ep [ \max_{1 \leq j \leq p} {\textstyle \sum}_{i=1}^{n} Z_{ij}^{2}  ] )^{1/2} \sqrt{\log p}. \quad (\text{Jensen})
\end{align*}
By Lemma \ref{lem: maximal ineq nonnegative} applied to $V_{ij} = Z_{ij}^{2}$, we have
\begin{equation*}
\Ep [ \max_{1 \leq j \leq p} {\textstyle \sum}_{i=1}^{n} Z_{ij}^{2}  ] \lesssim \sigma^{2} + \Ep [ M^{2} ] \log p.
\end{equation*}
This implies the desired conclusion.
\end{proof}

\section*{Acknowledgments}
V. Chernozhukov and D. Chetverikov  are supported by a National Science Foundation grant. K. Kato is supported by the Grant-in-Aid for Young Scientists (B) (22730179, 25780152), the Japan Society for the Promotion of Science. We would like to thank the Editors and an anonymous referee for their careful review.

\begin{thebibliography}{99}
\bibitem[Adler and Taylor(2007)]{AdlerTaylor2007}
Adler, R. and Taylor, J. (2007). {\em Random Fields and Geometry}. Springer.
\bibitem[Ball(1993)]{Ball1993}
Ball, K. (1993). The reverse isoperimetric problem for Gaussian measure. {\em Discrete Comput. Geom.} \textbf{10} 411-420.
\bibitem[Bentkus(2003)]{Bentkus2003}
Bentkus, V. (2003). On the dependence of the Berry-Esseen bound on dimension. {\em J. Statist. Plann. Infer.} \textbf{113} 385-402.
\bibitem[Borell(1975)]{B75}
Borell, C. (1975). The Brunn-Minkowski inequality in Gauss space. {\em Invent. Math.} \textbf{30} 205-216.
\bibitem[B\"{u}hlmann and van de Geer(2011)]{BV11}
B\"{u}hlmann, P. and van de Geer, S. (2011). {\em Statistics for High-Dimensional Data: Methods,Theory and Applications}. Springer.
% \bibitem[Chatterjee(2005a)]{Chatterjee2005a}
% Chatterjee, S. (2005a). A simple invariance theorem. arXiv:math/0508213.
\bibitem[Candes and Tao(2007)]{CandesTao2007}
Candes, E. and Tao, T. (2007). The Dantzig selector: statistical estimation when $p$ is much larger than $n$. {\em Ann. Statist.} \textbf{35} 2313-2351.
\bibitem[Chatterjee(2005)]{Chatterjee2005b}
Chatterjee, S. (2005). An error bound in the Sudakov-Fernique inequality. arXiv:math/0510424.
% \bibitem[Chatterjee and Meckes(2008)]{ChatterjeeMeckes2008}
% Chatterjee, S. and Meckes, E. (2008). Multivariate normal approximation using exchangeable pairs. {\em ALEA Lat. Am. J. Probab. Math. Stat.}
\bibitem[Chen et al.(2011)]{ChenGoldsteinShao2011}
Chen, L., Goldstein, L. and Shao, Q.-M.  (2011). {\em Normal Approximation by Stein's Method}. Springer.
%\bibitem[Chernozhukov et al.(2012)]{CCK12a}
%Chernozhukov, V., Chetverikov, D. and Kato, K. (2012). Gaussian approximation of suprema of empirical processes.
\bibitem[Chernozhukov et al.(2013a)]{CCK12}
Chernozhukov, V., Chetverikov, D. and Kato, K. (2013a). Gaussian approximations and multiplier bootstrap for maxima of sums of high-dimensional random vectors.  {\em Ann. Statist}. \textbf{41} 2786-2819.
\bibitem[Chernozhukov et al.(2013b)]{CCK13}
Chernozhukov, V., Chetverikov, D. and Kato, K. (2013b). Anti-concentration and honest, adaptive confidence bands. arXiv:1303:7152.
\bibitem[Dudley(1999)]{D99}
Dudley, R.M. (1999). {\em Uniform Central Limit Theorems}. Cambridge University Press.
% \bibitem[Galambos(1977)]{G77}
% Galambos, J. (1977). {\em The Asymptotic Theory of Extreme Order Statistics}. John Wiley \& Sons.
\bibitem[Gordon(1985)]{G85}
Gordon, Y. (1985). Some inequalities for Gaussian processes and applications. {\em Israel J. Math. } \textbf{50} 265-289.
%\bibitem[Harsha et al.(2012)]{HKM2012}
%Harsha, P., Klivans, A., and Meka, R. (2012). An invariance principle for polytopes. {\em Journal of the ACM} \textbf{59}.
\bibitem[Kilvans et al.(2008)]{Klivans08}
Klivans, A.R., O'Donnell, R., and Servedio, R.A. (2008). Learning geometric concepts via Gaussian surface area. In: {\em Proc. 49th IEEE Symp. on Foundations of Comp. Science (FOCS)}  pp. 541-550.
\bibitem[Leadbetter et al.(1983)]{Leadbetter1983}
Leadbetter, M., Lindgren, G. and Rootz\'{e}n, H.  (1983). {\em Extremes and Related Properties of Random Sequences and Processes}. Springer.
\bibitem[Ledoux(2001)]{L01}
Ledoux, M. (2001). {\em Concentration of Measure Phenomenon}. American Mathematical Society.
\bibitem[Ledoux and Talagrand(1991)]{LT91}
Ledoux, M. and Talagrand, M. (1991). {\em Probability in Banach Spaces}. Springer.
\bibitem[Li and Shao(1991)]{LS01}
Li, W. and Shao, Q.-M. (2001). Gaussian processes: inequalities, small ball probabilities and applications. In: {\em Handbook of Statistics Vol. 19},  pp. 533-597. North-Holland.
\bibitem[Li and Shao(2002)]{LS02}
Li, W. and Shao, Q.-M. (2002). A normal comparison inequality and applications. {\em Probab. Theory Relat. Fields} \textbf{122} 494-508.
\bibitem[Nazarov(2003)]{Nazarov03}
Nazarov, F. (2003). On the maximal perimeter of a convex set in $\RR^n$ with respect to a Gaussian measure. In: {\em Geometric Aspects of Functional Analysis}, Lecture Notes in Mathematics \textbf{1807}, pp. 169-187. Springer.
\bibitem[Nourdin and Viens(2009)]{NV09}
Nourdin, I. and Viens, F.G. (2009). Density formula and concentration inequalities with Malliavin calculus. {\em Electron. J. Probab.} \textbf{14} 2287-2309.
\bibitem[Panchenko(2013)]{Panchenko2013}
Panchenko, D. (2013). {\em The Sherrington-Kirkpatrick Model}. Springer.
\bibitem[Rudelson and Vershynin(2008)]{RV08}
Rudelson, M. and Vershynin, R. (2008).  The Littlewood-Offord problem and invertibility of random matrices. {\em Adv. Math.}  \textbf{218} 600-633.
\bibitem[Rudelson and Vershynin(2009)]{RV09}
Rudelson, M. and Vershynin, R. (2009). Smallest singular value of a random rectangular matrix. {\em Commun. Pure Appl. Anal.} \textbf{62} 1707-1739.
% \bibitem[Talagrand(1996)]{T96}
% Talagrand, M. (1996). New concentration inequalities in product spaces. {\em Invent. Math. } \textbf{126}  503-563.
\bibitem[Slepian(1962)]{Slepian1962}
Slepian, D. (1962). The one-sided barrier problem for Gaussian noise. {\em Bell Syst. Tech. J.} \textbf{41} 463-501.
% \bibitem[Smirnov(1950)]{Smirnov1950}
% Smirnov, N. (1950). On the construction of confidence regions for the density of distribution of random variables. {\em Doklady Akad. Nauk SSSR}  \textbf{74} 189-191.
\bibitem[Stein(1981)]{Stein1981}
Stein, C. (1981). Estimation of the mean of a multivariate normal distribution. {\em Ann. Statist.} \textbf{9} 1135-1151.
\bibitem[Sudakov and Tsirel'son(1978)]{ST78}
Sudakov, V.N. and Tsirel'son, B.S. (1978). Extremal properties of half-spaces for spherically invariant measures. {\em J. Soviet Math.} \textbf{9} 9-18.
\bibitem[Talagrand(2003)]{Talagrand2003}
Talagrand, M. (2003). {\em Spin Glasses: A Challenge for Mathematicians}.  Springer.
\bibitem[van der Vaart and Wellner(1996)]{VW96}
van der Vaart, A.W. and Wellner, J.A. (1996). {\em Weak Convergence and Empirical Processes: With Applications to Statistics}. Springer.
\bibitem[Vershynin and Rudelson(2007)]{VR07}
Vershynin, R. and Rudelson, M. (2007). Anti-concentration inequalities. Conference Proceeding for Phenomena in High Dimensions, Samos, Greece.
\bibitem[Yan(2009)]{Y09}
Yan, L. (2009). Comparison inequalities for one sided normal probabilities. {\em J. Theoret. Probab.} \textbf{22} 827-836.
\bibitem[Ylvisaker(1965)]{Y65}
Ylvisaker, D. (1965). The expected number of zeros of a stationary Gaussian process. {\em Ann. Math. Statist.} \textbf{36} 1043-1046.
\bibitem[Ylvisaker(1968)]{Y68}
Ylvisaker, D. (1968). A note on the absence of tangencies in Gaussian sample paths. {\em Ann. Math. Statist.} \textbf{39} 261-262.
\end{thebibliography}
\end{document}
